%% \SetAltSpoken: the connector \altn and \altg put between alternatives in
%% their spoken form.  It was hardcoded English until now -- "A, B, or C"
%% with no way to say "A, B ou C" -- which made a French or German document
%% announce an English word in the middle of its own alternatives.
%%
%% Three things vary and all three are asserted here: the word, the
%% punctuation between the earlier items, and whether that punctuation is
%% repeated before the connector (English does; French and German do not).
%% The default must come through UNCHANGED, and the starred form must be
%% able to reproduce it -- a hook that cannot express the default it
%% replaces hides one of the two behaviours from the author.
%%
%% Tagged preamble, so the /Alt strings are in the PDF.
\DocumentMetadata{uncompress,lang=en,testphase={phase-III}}
\documentclass{article}
\usepackage[lazy,gb4e]{linguexx}
\begin{document}


\ex. ALTDEFTWO \altn{aa}{bb} end.
\ex. ALTDEFTHREE \altn{cc}{dd}{ee} end.

%% The setting is global, so each variant is fenced in its own group; the
%% last one must not leak into the next.
{\SetAltSpoken{ou}
 \ex. ALTFRTWO \altn{ff}{gg} end.
 \ex. ALTFRTHREE \altn{hh}{ii}{jj} end.}

{\SetAltSpoken{oder}
 \ex. ALTDETHREE \altn{kk}{ll}{mm} end.}

%% The starred form asks for the serial comma back: this is the default
%% spelled out, and must match the default above exactly.
{\SetAltSpoken*{or}
 \ex. ALTSTARRED \altn{nn}{oo}{pp} end.}

%% The optional argument replaces the punctuation itself.
{\SetAltSpoken{y}[;]
 \ex. ALTPUNCT \altn{qq}{rr}{ss} end.}

%% \altg shares the builder with \altn, so the setting must reach a gloss
%% too -- and on BOTH tiers, since the stack is written twice.
{\SetAltSpoken{ou}
 \ex. \gll ALTGLOSS \altg{tt}{uu}\\
      objectword \altg{vv}{ww}\\
 \glt `a gloss whose stacks are spoken in French'}

%% Spacing is the package's business, so a word written with spaces around
%% it is the SAME setting -- otherwise the stored " ou " speaks as two
%% spaces a side, which shows up neither on the page nor in the log.
{\SetAltSpoken{ ou }
 \ex. ALTPADDED \altn{c1}{d1} end.}

%% The optional argument is trimmed for the same reason as the word.
{\SetAltSpoken{ y }[ ; ]
 \ex. ALTPADPUNCT \altn{e1}{f1}{g1} end.}

%% An empty connector is not an error: it collapses to a SINGLE space
%% rather than the two the naive spacing would give, which is how a list
%% joined by punctuation alone is written.
{\SetAltSpoken{}
 \ex. ALTEMPTY \altn{a1}{b1} end.}

%% Back outside every group: the default is still in force.
\ex. ALTAFTER \altn{xx}{yy}{zz} end.
\end{document}
